\section*{Capítulo \ref{ch:b1:biosyntheses-integration} --- Biossínteses e a célula integrada}

\begin{solution}{exo:b1:biosyntheses-integration:1}
Esqueletos de carbono (vindos das vias centrais), poder redutor (\omterm{def:b1:biosyntheses-integration:ppp}{NADPH})
e energia (\omterm{def:b1:water-small-molecules:atp}{ATP}). A \omterm{def:b1:biosyntheses-integration:ppp}{via das pentoses-fosfato}.
\end{solution}

\begin{solution}{exo:b1:biosyntheses-integration:2}
A hexoquinase, contornada pela glicose-6-fosfatase; a fosfofrutoquinase,
pela frutose-1,6-bisfosfatase; a piruvato-quinase, pela
piruvato-carboxilase mais a PEP-carboxiquinase.
\end{solution}

\begin{solution}{exo:b1:biosyntheses-integration:3}
Pico de cerca de \qty{7.5}{mmol/L} uma hora depois da refeição; de volta
a \qty{5}{} depois de cerca de três horas; mínimo de cerca de
\qty{4.5}{mmol/L} antes do café da manhã.
\end{solution}

\begin{solution}{exo:b1:biosyntheses-integration:4}
A semente tem o ciclo do glioxilato, que transforma dois \omterm{def:b1:respiration-fermentation:pdh}{acetil-CoA} num
ácido de quatro carbonos sem perder carbono como $\mathrm{CO_2}$, e a
partir desse ácido fabrica açúcar. Aos animais faltam as duas \omterm{def:b1:enzymes:enzyme}{enzimas} do
atalho: o \omterm{def:b1:respiration-fermentation:pdh}{acetil-CoA} deles entra no \omterm{def:b1:respiration-fermentation:krebs}{ciclo de Krebs} e os dois carbonos
dele saem como $\mathrm{CO_2}$ antes que apareça qualquer oxaloacetato
líquido.
\end{solution}

\begin{solution}{exo:b1:biosyntheses-integration:5}
Seis equivalentes de \omterm{def:b1:water-small-molecules:atp}{ATP} (duas carboxilações, dois GTP, dois \omterm{def:b1:water-small-molecules:atp}{ATP} na
etapa do fosfoglicerato) por glicose. Líquido por ciclo:
$2 - 6 = -4$ \omterm{def:b1:water-small-molecules:atp}{ATP}. Quem paga é o fígado, com oxigênio, depois da
disparada --- o músculo tomou emprestado um \omterm{def:b1:water-small-molecules:atp}{ATP} que não conseguia
fabricar a tempo em aerobiose.
\end{solution}

\begin{solution}{exo:b1:biosyntheses-integration:6}
Cada glicose dá 2 \omterm{def:b1:respiration-fermentation:pdh}{acetil-CoA}: 4 glicoses. Cada glicose que passa pelo
ramo oxidativo dá 2 \omterm{def:b1:biosyntheses-integration:ppp}{NADPH}: 7 glicoses. Onze glicoses para um palmitato
--- quatro fornecem carbono e sete fornecem elétrons.
\end{solution}

\begin{solution}{exo:b1:biosyntheses-integration:7}
\omterm{def:b1:enzymes:enzyme}{Enzimas} diferentes permitem que cada sentido seja ligado ou desligado de
forma independente; compartimentos diferentes (o \omterm{def:b1:eukaryotic-cell:organelle}{citosol} para a síntese,
a matriz para a oxidação) separam os dois conjuntos de intermediários; o
\omterm{def:b1:biosyntheses-integration:ppp}{NADPH} é mantido reduzido para a síntese enquanto o \omterm{def:b1:respiration-fermentation:carriers}{NAD} é mantido oxidado
para a degradação, de modo que cada sentido é termodinamicamente
favorecido no lugar dele. O malonil-CoA, o primeiro intermediário
comprometido com a síntese, inibe a lançadeira da carnitina que carrega
os \omterm{def:b1:lipids:fattyacid}{ácidos graxos} para dentro da \omterm{def:b1:eukaryotic-cell:mitochondrion}{mitocôndria}: enquanto \omterm{def:b1:lipids:triglyceride}{gordura} está sendo
fabricada, nenhuma é queimada.
\end{solution}

\begin{solution}{exo:b1:biosyntheses-integration:8}
Alanina $+ \alpha$-cetoglutarato $\rightleftharpoons$ piruvato $+$
\omterm{def:b1:biosyntheses-integration:nitrogen}{glutamato}. Depois de uma refeição rica em \omterm{def:b1:proteins:peptide}{proteína}, a alanina excedente
entrega o nitrogênio dela ao \omterm{def:b1:biosyntheses-integration:nitrogen}{glutamato} (para a direita), para o descarte
como ureia, e o carbono dela vira combustível ou glicose. Num jejum, a
\omterm{def:b1:proteins:peptide}{proteína} muscular é degradada e os grupos amino dela são recolhidos no
piruvato como alanina (para a esquerda, no músculo), enviados ao fígado,
e ali reconvertidos em piruvato e em glicose.
\end{solution}

\begin{solution}{exo:b1:biosyntheses-integration:9}
O fígado não consegue liberar glicose livre: a glicose do sangue cai
perigosamente entre as refeições (hipoglicemia). A glicose-6-fosfato se
acumula e é desviada para o \omterm{def:b1:carbohydrates:polysaccharide}{glicogênio}, que o fígado não consegue
mobilizar para o sangue: o fígado aumenta de tamanho, cheio de
\omterm{def:b1:carbohydrates:polysaccharide}{glicogênio}. O excesso de glicose-6-fosfato também desce a \omterm{def:b1:respiration-fermentation:glycolysis}{glicólise} até
o lactato, que sobe no sangue.
\end{solution}

\begin{solution}{exo:b1:biosyntheses-integration:10}
$120/0.55 = \qty{218}{g}$ de \omterm{def:b1:proteins:peptide}{proteína} por dia, cerca de \qty{1.1}{kg} de
músculo. Com corpos cetônicos, $40/0.55 = \qty{73}{g}$ de \omterm{def:b1:proteins:peptide}{proteína},
\qty{360}{g} de músculo: a perda cai em dois terços, e a sobrevivência
com a reserva de \omterm{def:b1:lipids:triglyceride}{gordura} passa de semanas a meses.
\end{solution}

\begin{solution}{exo:b1:biosyntheses-integration:11}
A mesma cascata de fosforilação ativa uma e inativa a outra, de modo que
o sinal que abre uma válvula fecha a outra. Se as duas estivessem
ativas, cada glicose acrescentada (um UTP, equivalente a um \omterm{def:b1:water-small-molecules:atp}{ATP}) e
retirada (de graça) custaria um \omterm{def:b1:water-small-molecules:atp}{ATP} por volta e não produziria nada além
de calor. A mamangava roda exatamente um ciclo desses (sobre a
frutose-6-fosfato) nos músculos de voo dela antes da decolagem numa
manhã fria, queimando \omterm{def:b1:water-small-molecules:atp}{ATP} como aquecedor até que o músculo esteja quente
o bastante para voar.
\end{solution}

\begin{solution}{exo:b1:biosyntheses-integration:12}
Cada \omterm{def:b1:enzymes:enzyme}{enzima} de encruzilhada responde apenas às moléculas que a cercam
--- \omterm{def:b1:water-small-molecules:atp}{ATP}, AMP, citrato, \omterm{def:b1:respiration-fermentation:pdh}{acetil-CoA}, um fosfato colocado nela por uma
quinase --- e abre ou fecha conforme isso; \omterm{def:b1:enzymes:enzyme}{enzimas} separadas para os
dois sentidos fazem com que abrir um sentido feche o outro. Nenhuma
\omterm{def:b1:enzymes:enzyme}{enzima} sabe do café da manhã nem do encéfalo; e no entanto a soma dessas
respostas locais é que o fígado armazena depois de uma refeição, libera
à noite e poupa \omterm{def:b1:proteins:peptide}{proteína} num jejum: uma estratégia coerente, sem
estrategista. O ``plano'' é a ligação das válvulas, escolhida pela
seleção, e os hormônios que ajustam muitas válvulas de uma vez são o que
faz o \omterm{def:b1:organism-environment:organism}{organismo} agir como um só.
\end{solution}

\begin{solution}{pb:b1:biosyntheses-integration:1}
\textbf{1.} \omterm{def:b1:carbohydrates:polysaccharide}{Glicogênio} hepático $100\times 17 = \qty{1.7}{MJ}$, 0.2 dia;
\omterm{def:b1:carbohydrates:polysaccharide}{glicogênio} muscular \qty{6.8}{MJ}, 0.8 dia; \omterm{def:b1:lipids:triglyceride}{gordura} $12\,000\times 38 =
\qty{456}{MJ}$, 54 dias; \omterm{def:b1:proteins:peptide}{proteína} \qty{170}{MJ}, 20 dias (nunca
plenamente utilizável).
\textbf{2.} O \omterm{def:b1:carbohydrates:polysaccharide}{glicogênio} hepático: o fígado tem glicose-6-fosfatase. O
músculo não a tem, de modo que o \omterm{def:b1:carbohydrates:polysaccharide}{glicogênio} dele só rende
glicose-6-fosfato para a própria \omterm{def:b1:respiration-fermentation:glycolysis}{glicólise}.
\textbf{3.} $120 + 40 = \qty{160}{g}$ por dia.
\textbf{4.} $0.005\times 5\times 180 = \qty{4.5}{g}$: a \qty{5}{g/h}, cerca de \qty{54}{min}.
\textbf{5.} $120\times 17 = \qty{2.0}{MJ}$, um quarto dos \qty{8.4}{MJ}.
\textbf{6.} A conversão em \omterm{def:b1:lipids:fattyacid}{ácidos graxos} e em \omterm{def:b1:lipids:triglyceride}{triglicerídeo} no fígado,
exportados para o \omterm{def:b1:body-plans-tissues:tissue}{tecido} adiposo; a insulina.
\textbf{7.} $160/24 = \qty{6.7}{g/h}$.
\textbf{8.} $100/6.7 = \qty{15}{h}$.
\textbf{9.} O \omterm{def:b1:carbohydrates:polysaccharide}{glicogênio} fornece dois terços, \qty{4.4}{g/h}: $100/4.4 =
\qty{22}{h}$.
\textbf{10.} Às 11:00 da manhã seguinte pela primeira estimativa, e às
18:00 pela segunda --- pela noite adentro nos dois casos, com a margem
dependendo de quão cedo a \omterm{def:b1:biosyntheses-integration:gluconeogenesis}{gliconeogênese} assume.
\textbf{11.} Gasto de repouso ao longo de \qty{12}{h}: \qty{4.2}{MJ};
\qty{60}{\%} disso são \qty{2.5}{MJ}; $2500/38 = \qty{66}{g}$ de
\omterm{def:b1:lipids:triglyceride}{gordura}.
\textbf{12.} $160/0.55 = \qty{290}{g}$ de \omterm{def:b1:proteins:peptide}{proteína}, ou seja,
\qty{1.45}{kg} de músculo por dia.
\textbf{13.} $160/180 = \qty{0.89}{mol}$ de glicose; a 6 \omterm{def:b1:water-small-molecules:atp}{ATP} cada uma,
são \qty{5.3}{mol} de \omterm{def:b1:water-small-molecules:atp}{ATP}, e a \qty{50}{kJ/mol}, cerca de \qty{265}{kJ}
--- uns \qty{3}{\%} dos \qty{8.4}{MJ} do dia: fabricar a glicose é
barato, e o que falta ao corpo são os esqueletos de carbono, e não a
energia.
\textbf{14.} Gasto não glicídico $8.4 - 160\times 0.017 =
\qty{5.7}{MJ}$: $5700/38 = \qty{150}{g}$ de \omterm{def:b1:lipids:triglyceride}{gordura}, cujo glicerol dá
\qty{15}{g} de glicose.
\textbf{15.} $(160 - 15)/0.55 = \qty{264}{g}$ de \omterm{def:b1:proteins:peptide}{proteína},
\qty{1.3}{kg} de músculo por dia.
\textbf{16.} \qty{3}{kg} de \omterm{def:b1:proteins:peptide}{proteína} muscular a \qty{264}{g} por dia:
cerca de \qty{11}{d}.
\textbf{17.} $(80 - 15)/0.55 = \qty{118}{g}$ de \omterm{def:b1:proteins:peptide}{proteína}, \qty{0.6}{kg}
de músculo por dia; metade do músculo em \qty{25}{d} --- e na prática a
perda cai ainda mais, para \qtyrange{20}{30}{g} por dia, à medida que o
corpo economiza.
\textbf{18.} $456/8.4 = 54$ dias; a \qty{6.5}{MJ}, 70 dias.
\textbf{19.} A \omterm{def:b1:proteins:peptide}{proteína} não é uma reserva, e sim a maquinaria: perder um
terço dela (músculos respiratórios, coração, \omterm{def:b1:proteins:peptide}{proteínas} imunes) é fatal,
e a \qty{100}{g} por dia ou mais isso acontece em poucas semanas,
enquanto a \omterm{def:b1:lipids:triglyceride}{gordura} duraria dois meses. Os corpos cetônicos permitem que
o encéfalo funcione com \omterm{def:b1:lipids:triglyceride}{gordura}, cortam a demanda de glicose e, com
isso, a perda de \omterm{def:b1:proteins:peptide}{proteína} a um terço, e estendem a sobrevivência até o
limite fixado pela \omterm{def:b1:lipids:triglyceride}{gordura}.
\textbf{20.} A glicogênio-fosforilase ativa às 03:00, e a
glicogênio-sintase às 09:00.
\textbf{21.} O glucagon (03:00) e a insulina (09:00); a fosforilação das
duas \omterm{def:b1:enzymes:enzyme}{enzimas} por uma cascata de quinases, revertida por uma fosfatase.
\textbf{22.} O \omterm{def:b1:respiration-fermentation:pdh}{acetil-CoA} ativa a piruvato-carboxilase; o \omterm{def:b1:water-small-molecules:atp}{ATP} (e o
citrato) inibe a fosfofrutoquinase.
\textbf{23.} Glicose a piruvato rende 2 \omterm{def:b1:water-small-molecules:atp}{ATP}; piruvato de volta a glicose
custa 6: cada ida e volta queima 4 \omterm{def:b1:water-small-molecules:atp}{ATP} sem produto nenhum. Com \omterm{def:b1:enzymes:enzyme}{enzimas}
distintas nos três desvios, os mesmos sinais (\omterm{def:b1:water-small-molecules:atp}{ATP}, \omterm{def:b1:respiration-fermentation:pdh}{acetil-CoA},
fosforilação) ativam um conjunto e inibem o outro, de modo que só um
sentido fica aberto.
\textbf{24.} Sem insulina as válvulas ficam ajustadas como num jejum,
seja qual for a glicose: a fosforilase, a \omterm{def:b1:biosyntheses-integration:gluconeogenesis}{gliconeogênese} e a lipólise
estão ligadas, a síntese está desligada, e o fígado, ``acreditando'' que
o corpo passa fome, despeja glicose e cetonas num sangue já cheio de
glicose que os \omterm{def:b1:body-plans-tissues:tissue}{tecidos} dependentes de insulina não conseguem captar.
\textbf{25.} Cerca de \qty{15}{h} sem \omterm{def:b1:biosyntheses-integration:gluconeogenesis}{gliconeogênese} e \qty{22}{h} com
ela --- uma noite e uma manhã; perda de \omterm{def:b1:proteins:peptide}{proteína} de cerca de
\qty{260}{g} por dia no início e de \qty{120}{g} (caindo mais ainda)
assim que os corpos cetônicos alimentam o encéfalo; a reserva de \omterm{def:b1:lipids:triglyceride}{gordura}
dura uns \qtyrange{55}{70}{d}.
\end{solution}
